\section*{第 \ref{ch:g10:storms-and-cyclones} 章——风暴与热带气旋}

\begin{solution}{exo:g10:storms-and-cyclones:1}
\omterm{def:g10:storms-and-cyclones:thunderstorm}{雷暴}是由积雨云产生的风暴，伴有\omterm{def:g10:storms-and-cyclones:lightning}{闪电}和
雷声，通常伴有强降水和阵风，有时还有冰雹。\omterm{def:g10:storms-and-cyclones:lightning}{闪电}是一种巨大的
电火花，使风暴中的带电区域放电。雷声是
\omterm{def:g10:storms-and-cyclones:lightning}{闪电}通道中的空气骤然受热膨胀，
产生冲击波所发出的声音。
\end{solution}

\begin{solution}{exo:g10:storms-and-cyclones:2}
\omterm{def:g10:storms-and-cyclones:tropical-cyclone}{热带气旋}是一种直径数百千米的风暴，形成于
温暖的热带海洋上，围绕低压中心，气流螺旋式向内
汇聚，并伴有\omterm{def:g10:storms-and-cyclones:thunderstorm}{雷暴}带。\omterm{def:g10:storms-and-cyclones:tropical-cyclone}{风暴眼}是其平静、几乎无云的中心；
\omterm{def:g10:storms-and-cyclones:tropical-cyclone}{眼墙}是环绕\omterm{def:g10:storms-and-cyclones:tropical-cyclone}{风暴眼}的高耸\omterm{def:g10:storms-and-cyclones:thunderstorm}{雷暴}环。风力
在\omterm{def:g10:storms-and-cyclones:tropical-cyclone}{眼墙}处最强。
\end{solution}

\begin{solution}{exo:g10:storms-and-cyclones:3}
\num{70}~节：1 级。\num{120}~节：4 级。
\num{140}~节：5 级。最后两者为强飓风（3
级或以上）。
\end{solution}

\begin{solution}{exo:g10:storms-and-cyclones:4}
$343 \times 9 \approx \qty{3.1}{km}$，然后 $343 \times 15 \approx
\qty{5.1}{km}$。第二次\omterm{def:g10:storms-and-cyclones:lightning}{闪电}距离更远：风暴正在远离。
\end{solution}

\begin{solution}{exo:g10:storms-and-cyclones:5}
\omterm{def:g10:storms-and-cyclones:lightning}{闪电}往往击中周围最高的点：
孤立立于田野中的一棵树就是这样的点，\omterm{def:g10:storms-and-cyclones:lightning}{闪电}电流沿树干向下，
并在其周围的地面中扩散。应进入建筑物或封闭的汽车内，
并待到最后一声雷声过后很久再出来。
\end{solution}

\begin{solution}{exo:g10:storms-and-cyclones:6}
1 节约为 $1852/3600 \approx \qty{0.514}{m/s}$。\num{64}~节：$64 \times
0.514 \approx \qty{32.9}{m/s}$，$32.9 \times 3.6 \approx
\qty{119}{km/h}$。\num{137}~节：$\qty{70.5}{m/s}$，约
\qty{254}{km/h}（官方换算表经另一单位取整后给出
\qty{252}{km/h}）。
\end{solution}

\begin{solution}{exo:g10:storms-and-cyclones:7}
每天约 $44 \times 86\,400 \approx \num{3.8e6}$ 次\omterm{def:g10:storms-and-cyclones:lightning}{闪电}，全年 $3.8\times10^6
\times 365.25 \approx \num{1.4e9}$，即每年约 14 亿次。
\end{solution}

\begin{solution}{exo:g10:storms-and-cyclones:8}
$\Delta p = 1010 - 935 = \qty{75}{hPa}$：水面隆起约 \qty{75}{cm}。
\end{solution}

\begin{solution}{exo:g10:storms-and-cyclones:9}
它们形成于温暖的热带海洋上，多位于两半球的
\qty{10}{\degree}至\qty{30}{\degree}纬度之间，南大西洋
极少。它们先向西漂移，随后向两极方向偏转，
并在较冷的水域上减弱。赤道附近没有这类风暴，因为那里
地球自转使流入的空气无法转向；没有这一转向，空气
直接流入低压区并将其填满，旋转风暴因而无法形成。
\end{solution}

\begin{solution}{exo:g10:storms-and-cyclones:10}
$1031 - 1.56 \times 70 \approx \qty{922}{hPa}$；$1031 - 1.56 \times 40
\approx \qty{969}{hPa}$。
\end{solution}

\begin{solution}{exo:g10:storms-and-cyclones:11}
风暴登陆后失去蒸发水汽的供应，
也就失去了驱动它的\omterm{def:g10:earth-energy-budget:latent-heat}{潜热}；粗糙的地表还会减慢风速。在寒冷
的水面上，蒸发要弱得多，因为冷空气在\omterm{def:g10:why-clouds-form:humidity}{饱和}时
所能容纳的水汽少得多，\omterm{def:g10:earth-energy-budget:latent-heat}{潜热}的供应随之减少。
\end{solution}

\begin{solution}{exo:g10:storms-and-cyclones:12}
在 \qty{0}{\degreeCelsius} 时：$\sqrt{1.4 \times 8.314 \times 273.15/0.02897}
\approx \qty{331}{m/s}$；在 \qty{30}{\degreeCelsius} 时，取 $T =
\qty{303.15}{K}$，声速约为 \qty{349}{m/s}。对于 \qty{9}{s}：$9 \times 331 \approx
\qty{2980}{m}$，$9 \times 349 \approx \qty{3140}{m}$，两者相差
约 \qty{160}{m}，约占 \qty{5}{\percent}。该经验规则给出 \qty{3}{km}：
用来判断风暴远近已经足够，这正是它的用途。
\end{solution}

\begin{solution}{exo:g10:storms-and-cyclones:13}
$\Delta p = 1010 - 940 = \qty{70}{hPa}$：隆起约 \qty{0.7}{m}，
风的作用部分为 $4.0 - 0.7 = \qty{3.3}{m}$，是前者四倍多。高潮时
水位高出平均海平面 $4.0 + 1.5 = \qty{5.5}{m}$。
\end{solution}

\begin{solution}{exo:g10:storms-and-cyclones:14}
\omterm{def:g10:storms-and-cyclones:tornado}{龙卷风}：一圈长 $2\pi \times 100 \approx \qty{628}{m}$，
需时 $628/90 \approx \qty{7}{s}$。飓风：$2\pi \times 30\,000 \approx
\qty{188}{km}$，需时 $188\,000/60 \approx \qty{3140}{s}$，约
\qty{52}{min}。\omterm{def:g10:storms-and-cyclones:tornado}{龙卷风}的转速约为飓风的 450 倍：前者是小型、
极快的旋风，后者是巨大而较慢的旋涡。
\end{solution}

\begin{solution}{exo:g10:storms-and-cyclones:15}
$24/3888 \approx \qty{0.6}{\percent}$；$414/3888 \approx
\qty{10.6}{\percent}$；$121/3888 \approx \qty{3.1}{\percent}$。若纬度
不起作用，约 \qty{9}{\percent} 的风暴（约 350 个）会
起源于赤道 \qty{5}{\degree} 以内；实际只有 24 个，少约十五倍。
赤道附近地球自转太弱，难以启动旋转。
\end{solution}

\begin{solution}{pb:g10:storms-and-cyclones:1}
\textbf{1.} $\pi \times (665\,000)^2 \approx \qty{1.39e12}{m^2}$.

\textbf{2.} $0.015 \times 1.39\times10^{12} \approx \qty{2.08e10}{m^3}$.

\textbf{3.} 一立方千米等于 \qty{e9}{m^3}：约 \qty{21}{km^3}。

\textbf{4.} $2.08\times10^{10} \times 1000 \approx \qty{2.08e13}{kg}$.

\textbf{5.} $2.08\times10^{13} \times 2.5\times10^6 \approx \qty{5.2e19}{J}$.

\textbf{6.} $5.2\times10^{19}/86\,400 \approx \qty{6.0e14}{W}$.

\textbf{7.} $6.0\times10^{14}/1.22\times10^{17} \approx 0.005$，约
\qty{0.5}{\percent}。

\textbf{8.} $240 \times 1.39\times10^{12} \approx \qty{3.3e14}{W}$。风暴
释放的热量几乎是其整个区域所吸收阳光热量的两倍。

\textbf{9.} 风暴并非依靠照射在其上的阳光运转：它
释放的是以水汽形式储存的热量，这些水汽由更大范围的暖
海面蒸发而来，并被风卷入。风暴把在大得多的区域内汇集的能量
集中释放。

\textbf{10.} 与降雨的质量相同，约每天 \qty{2.1e13}{kg}（
流入的水汽须与降雨流失的速度同步得到补充）。

\textbf{11.} $6.0\times10^{14}/82 \approx \qty{7.3e12}{m^2}$，约
为风暴面积的五倍：风暴必须从广得多的区域汲取水汽，并且
其强风使其下方海面的蒸发速度远高于平均水平。

\textbf{12.} 在 \qty{28}{\degreeCelsius} 时为 \qty{24.1}{g/kg}，
而在 \qty{20}{\degreeCelsius} 时为 \qty{14.7}{g/kg}：$24.1/14.7 \approx 1.64$，约
多出 \qty{64}{\percent}。

\textbf{13.} 暖水蒸发快，暖空气所能容纳的水汽多得多：
流入的空气携带的待释放\omterm{def:g10:earth-energy-budget:latent-heat}{潜热}多得多，并且，空气\omterm{def:g10:why-clouds-form:humidity}{饱和}后，
便可自由上升。在较冷的海面上，供给量太小，不足以驱动
热机。

\textbf{14.} 在陆地上没有可供蒸发的海面：\omterm{def:g10:earth-energy-budget:latent-heat}{潜热}的供给
中断，粗糙的地面还会使风减弱。

\textbf{15.} $1.5\times10^{12} \times 86\,400 \approx \qty{1.3e17}{J}$.

\textbf{16.} $96 \times 1852/3600 \approx \qty{49.4}{m/s}$.

\textbf{17.} $\frac12 \times 1.16 \times 49.4^2 \approx \qty{1.4e3}{J}$.

\textbf{18.} $5.2\times10^{19}/1.3\times10^{17} = 400$：\omterm{def:g10:earth-energy-budget:latent-heat}{潜热}
约为交给风的能量的 400 倍。

\textbf{19.} 其中大部分加热了风暴上层的空气，这些空气在顶部向外扩散，
并以红外辐射的形式将热量散失到太空中；风本身则
因克服海面和地面的摩擦力而损失能量。

\textbf{20.} $1/400 = \qty{0.25}{\percent}$（脚本给出的比值为
402）：\textbf{凝结释放的热量中，约每 400 份
有 1 份转化为风}。
\end{solution}
